% GENERATED FILE. Edit canonical lessons, then run npm run generate:books.
% canonical-source-hash: fnv1a64:98bbc9b98b38baad
% canonical-lessons: IT-C522-scolpire, IT-C522-scultore, IT-C522-marmo, IT-C522-argilla, IT-C522-modellare

\chapter{Marmo e argilla}
\label{ch:it-marmo-e-argilla}
\hlchaptermodality{522}

\begin{chapteropening}
\textbf{By the end of this chapter:} \emph{I can talk about sculpture and working with clay.}

Complete the last lesson of chapter 522: I can talk about sculpture and working with clay.
\end{chapteropening}

\section[scolpire]{scolpire — to sculpt, to carve}

\label{lesson:IT-C522-scolpire}



\begin{quote}
Before the new one: say the Italian for easel, then the Italian for watercolour.
\end{quote}

\subsection*{You'll want to know: scolpire}
\textbf{scolpire} — \textquotedblleft{}to sculpt, to carve\textquotedblright{}.

\begin{quote}
Ha scolpito una piccola statua nel legno.
\end{quote}

\begin{quote}
\emph{He carved a small statue out of wood.}
\end{quote}

An \emph{-isc-} verb: \emph{scolpisco}.

\begin{cousinweb}
\emph{scolpire} comes from Latin \emph{sculpere} \textquotedblleft{}to carve, to chisel\textquotedblright{} (EtymDB).
\end{cousinweb}

\begin{grammarlens}[title={What you've built}]
One more word for telling, explaining and describing at B1.
\end{grammarlens}

\subsection*{Your turn}
\begin{itemize}
\raggedright
  \item \emph{Say it:} \emph{scolpire}
  \item \emph{Say it:} \emph{scolpire}, once more
  \item \emph{Say it:} a sentence of your own with \emph{scolpire}
  \item \emph{Recall:} say the Italian for easel, then the Italian for watercolour, then say \emph{scolpire} again
\end{itemize}

\begin{tcolorbox}[breakable,colback=teal!4,colframe=teal!35!black,title={Before you move on}]
What does scolpire mean? (\textquotedblleft{}To sculpt, to carve\textquotedblright{}.) Say it once more.
\end{tcolorbox}
\section[lo scultore]{lo scultore — sculptor}

\label{lesson:IT-C522-scultore}



\begin{quote}
Before the new one: say the Italian for watercolour, then the Italian for to sculpt.
\end{quote}

\subsection*{You'll want to know: lo scultore}
\textbf{lo scultore} — \textquotedblleft{}sculptor\textquotedblright{}.

\begin{quote}
Lo scultore lavora nel suo studio.
\end{quote}

\begin{quote}
\emph{The sculptor works in his studio.}
\end{quote}

A woman is \emph{la scultrice}.

\begin{cousinweb}
\emph{scultore} is a learned borrowing from Latin \emph{sculptor} (EtymDB), from \emph{sculpere} \textquotedblleft{}to carve\textquotedblright{}.
\end{cousinweb}

\begin{grammarlens}[title={What you've built}]
One more word for telling, explaining and describing at B1.
\end{grammarlens}

\subsection*{Your turn}
\begin{itemize}
\raggedright
  \item \emph{Say it:} \emph{lo scultore}
  \item \emph{Say it:} \emph{lo scultore}, once more
  \item \emph{Say it:} a sentence of your own with \emph{lo scultore}
  \item \emph{Recall:} say the Italian for watercolour, then the Italian for to sculpt, then say \emph{lo scultore} again
\end{itemize}

\begin{tcolorbox}[breakable,colback=teal!4,colframe=teal!35!black,title={Before you move on}]
What does lo scultore mean? (\textquotedblleft{}Sculptor\textquotedblright{}.) Say it once more.
\end{tcolorbox}
\section[il marmo]{il marmo — marble}

\label{lesson:IT-C522-marmo}



\begin{quote}
Before the new one: say the Italian for to sculpt, then the Italian for sculptor.
\end{quote}

\subsection*{You'll want to know: il marmo}
\textbf{il marmo} — \textquotedblleft{}marble\textquotedblright{}.

\begin{quote}
In questa chiesa tutto è di marmo bianco.
\end{quote}

\begin{quote}
\emph{In this church everything is white marble.}
\end{quote}

\emph{Freddo come il marmo} — cold as marble.

\begin{cousinweb}
\emph{marmo} is inherited from Latin \emph{marmor} \textquotedblleft{}marble\textquotedblright{} (EtymDB), from Greek \emph{mármaros} \textquotedblleft{}a shining stone\textquotedblright{}.
\end{cousinweb}

\begin{grammarlens}[title={What you've built}]
One more word for telling, explaining and describing at B1.
\end{grammarlens}

\subsection*{Your turn}
\begin{itemize}
\raggedright
  \item \emph{Say it:} \emph{il marmo}
  \item \emph{Say it:} \emph{il marmo}, once more
  \item \emph{Say it:} a sentence of your own with \emph{il marmo}
  \item \emph{Recall:} say the Italian for to sculpt, then the Italian for sculptor, then say \emph{il marmo} again
\end{itemize}

\begin{tcolorbox}[breakable,colback=teal!4,colframe=teal!35!black,title={Before you move on}]
What does il marmo mean? (\textquotedblleft{}Marble\textquotedblright{}.) Say it once more.
\end{tcolorbox}
\section[l'argilla]{l'argilla — clay}

\label{lesson:IT-C522-argilla}



\begin{quote}
Before the new one: say the Italian for sculptor, then the Italian for marble.
\end{quote}

\subsection*{You'll want to know: l'argilla}
\textbf{l'argilla} — \textquotedblleft{}clay\textquotedblright{}.

\begin{quote}
I bambini fanno animali con l'argilla.
\end{quote}

\begin{quote}
\emph{The children make animals out of clay.}
\end{quote}

\emph{Un vaso di argilla} — a clay pot.

\begin{cousinweb}
\emph{argilla} is a learned borrowing from Latin \emph{argilla} \textquotedblleft{}white clay\textquotedblright{} (Lewis \& Short), from Greek \emph{árgillos}.
\end{cousinweb}

\begin{grammarlens}[title={What you've built}]
One more word for telling, explaining and describing at B1.
\end{grammarlens}

\subsection*{Your turn}
\begin{itemize}
\raggedright
  \item \emph{Say it:} \emph{l'argilla}
  \item \emph{Say it:} \emph{l'argilla}, once more
  \item \emph{Say it:} a sentence of your own with \emph{l'argilla}
  \item \emph{Recall:} say the Italian for sculptor, then the Italian for marble, then say \emph{l'argilla} again
\end{itemize}

\begin{tcolorbox}[breakable,colback=teal!4,colframe=teal!35!black,title={Before you move on}]
What does l'argilla mean? (\textquotedblleft{}Clay\textquotedblright{}.) Say it once more.
\end{tcolorbox}
\section[modellare]{modellare — to model, to shape}

\label{lesson:IT-C522-modellare}



\begin{quote}
Before the new one: say the Italian for marble, then the Italian for clay.
\end{quote}

\subsection*{You'll want to know: modellare}
\textbf{modellare} — \textquotedblleft{}to model, to shape\textquotedblright{}.

\begin{quote}
Modello l'argilla con le mani.
\end{quote}

\begin{quote}
\emph{I shape the clay with my hands.}
\end{quote}

Built on \emph{modello} \textquotedblleft{}model\textquotedblright{}.

\begin{cousinweb}
\emph{modellare} is built on \emph{modello} \textquotedblleft{}model\textquotedblright{}, from a little form of Latin \emph{modulus} \textquotedblleft{}a small measure\textquotedblright{} (Lewis \& Short).
\end{cousinweb}

\begin{grammarlens}[title={What you've built}]
One more word for telling, explaining and describing at B1.
\end{grammarlens}

\subsection*{Your turn}
\begin{itemize}
\raggedright
  \item \emph{Say it:} \emph{modellare}
  \item \emph{Say it:} \emph{modellare}, once more
  \item \emph{Say it:} a sentence of your own with \emph{modellare}
  \item \emph{Recall:} say the Italian for marble, then the Italian for clay, then say \emph{modellare} again
\end{itemize}

\begin{tcolorbox}[breakable,colback=teal!4,colframe=teal!35!black,title={Before you move on}]
What does modellare mean? (\textquotedblleft{}To model, to shape\textquotedblright{}.) Say it once more.
\end{tcolorbox}
